\subsection{Motivation}
Climate and environmental science depend on repeat optical imaging of regions such as tropical rainforests, where cloud cover is frequent and unpredictable at the timescale of a satellite pass.
Mission operators today often rely on deterministic strategies---fixed nadir attitude with pre-planned shutter times---that treat every target along a ground track equally.
When clouds occlude the line of sight, those strategies still consume onboard memory and downlink budget on low-value frames, while scientists lose observations that may not be recoverable on the next pass.
The cost is therefore threefold: reduced science yield, wasted capture capacity, and unnecessary actuator activity from indiscriminate pointing and slewing.

Autonomous onboard decision-making is motivated by latency: ground-in-the-loop replanning cannot react quickly enough as cloud fields evolve during overflight.
Both \emph{when} to expose the shutter and \emph{how} to orient the spacecraft for maximal image quality therefore matter, and they must respect attitude safety limits on reaction-wheel use and pointing envelopes.
This project treats clouds as the defining challenge rather than a secondary nuisance, and uses simulated vision---primary and secondary cameras---as the feedback channel for a learned ``smart autopilot'' that respects the same safety guidelines as the deterministic baseline.

For exposition we use a simplified meridian target stripe (many coordinate-defined targets over a rainforest corridor); the orbital geometry is not tied to a polar versus equatorial regime---any circular orbit viewed at the instant of overflight is equivalent for the control problem posed here.

\subsection{Objective}
\textbf{Primary objective.}
Demonstrate, in simulation, that a cloud-aware learned control policy can outperform the deterministic baseline (notebook~07 overflight) on observation value, defined as successful captures weighted by image quality, coverage, and target novelty, while cloud-blocked fractions reduce credit.

\textbf{Hypothesis.}
Reinforcement learning is appropriate because capture decisions are sequential and uncertain: cloud occlusion, attitude dynamics, and discrete shutter events interact over a single pass.
We test whether an MPO-trained policy can capitalize on vision feedback to avoid blind shuttering and prioritize high-value targets under a per-orbit capture budget---algorithm choice is secondary to this behavioral question.

\textbf{Secondary objectives.}
\begin{itemize}
  \item specify a reward aligned with mission constraints (capture quality, cloud occlusion, area novelty);
  \item evaluate against a documented baseline under \emph{matched} environment setups (same altitude draw, clouds, and target layout per comparison episode);
  \item produce reproducible run artifacts and demonstration exports (baseline versus trained-policy videos, learning curves, and comparison plots).
\end{itemize}

\textbf{Scope and simplifications.}
The model deliberately stays close to the control problem while omitting much operational detail: 2D attitude dynamics, basic circular-orbit geometry, physical ray tracing for cameras and clouds, safety arbitration, and single-pass episodes---no multi-orbit scheduling, flight hardware, or real imagery ingestion.
Section~\ref{sec:methods} lists modeled subsystems explicitly; Discussion states extensions (three-dimensional attitude, richer terrain, stronger RL benchmarks) as logical next steps rather than deliverables of this semester project.

\textbf{Long-term context.}
The intended architecture is hierarchical: a future mission planner would set observation intent while on-board control executes tracking.
Direct wheel-torque learning is the interim experimental interface used here to probe cloud-aware pointing behavior before that stack exists.

\subsection{Contributions}
This report documents:
\begin{itemize}
  \item \textbf{Platform.} A reproducible simulation-and-learning pipeline: canonical episode rollouts via \texttt{run\_simulation()} and \texttt{SimulationStateSeries}, seeded mission randomization (altitude and clouds), view-only rendering, notebook-driven training workflow with preflight gates, and fingerprinted warmup caching so baseline and learned policies replay identical stochastic missions.
  \item \textbf{Mission model.} Cloud-aware reward design and dual-camera observations for rainforest-style overflight under attitude safety constraints.
  \item \textbf{Evaluation.} Safety-constrained MPO training and matched baseline evaluation (notebook~07 overflight) under shared environment draws.
  \item \textbf{Evidence.} Qualitative and quantitative comparison of learned versus deterministic capture strategies from frozen run artifacts.
\end{itemize}

Substantial semester effort went into validated simulation and evaluation plumbing so that reinforcement-learning experiments could be compared credibly; initial learning results are reported alongside this engineering deliverable (Section~\ref{sec:discussion}).

\subsection{Report structure}
Section~\ref{sec:methods} describes the environment and learning method.
Section~\ref{sec:results} presents experiments.
Sections~\ref{sec:discussion} and~\ref{sec:conclusion} interpret limitations and outlook.
Technical constants and hyperparameters are maintained in \texttt{docs/presentation/} as the live record aligned with code.
